\documentclass[aspectratio=169,10pt,t]{beamer}
\usepackage{ac1}
\aulainfo{26}{Lançamento do Projeto Final}{Temas, grupos, escopo e cronograma}

\begin{document}
\capa

\begin{frame}{O projeto final}
  \begin{columns}[T]
    \column{0.55\textwidth}
    \begin{itemize}
      \item Programa \textbf{completo} que resolve um problema real da sua área de engenharia
      \item Em \textbf{grupos de 2 ou 3 alunos}
      \item Desenvolvido em \textbf{etapas}, com feedback contínuo em cada aula
      \item Integra tudo o que vimos: funções, listas e dicionários, arquivos, exceções, NumPy, Pandas e Matplotlib
    \end{itemize}
    \column{0.42\textwidth}
    \begin{exampleblock}{Objetivo}
      Passar de ``resolver exercícios'' para \alert{construir um programa} que outra pessoa consiga usar e entender.
    \end{exampleblock}
  \end{columns}
\end{frame}

\begin{frame}{O que deve aparecer no projeto}
  \begin{itemize}
    \item \textbf{Pense antes de programar} (Cap.~9): problema bem definido, esboço, casos-limite
    \item \textbf{Estruturas de controle} (Cap.~3): decisões e repetições
    \item \textbf{Funções} (Cap.~4): programa dividido em funções coesas, módulos organizados
    \item \textbf{Estruturas de dados} (Cap.~5): listas, dicionários, tuplas; matrizes
    \item \textbf{Arquivos} (Cap.~6): leitura e/ou gravação (texto, CSV)
    \item \textbf{Tratamento de erros}: entradas inválidas e arquivos ausentes
    \item \textbf{NumPy / Pandas / Matplotlib} quando aplicável: cálculo, análise e gráficos
  \end{itemize}
  \begin{alertblock}{Requisito mínimo}
    Pelo menos \textbf{5 funções}, \textbf{1 arquivo de dados}, \textbf{tratamento de erros} e, quando houver dados numéricos, \textbf{2 gráficos}.
  \end{alertblock}
\end{frame}

\begin{frame}{Sugestões de temas (gerais)}
  \begin{columns}[T]
    \column{0.5\textwidth}
    \begin{itemize}
      \item Sistema de biblioteca simples
      \item Controle financeiro pessoal
      \item Jogo textual (forca, velha, labirinto)
      \item Simulador simples: fila, elevador, semáforo
      \item Agenda de contatos / de tarefas
    \end{itemize}
    \column{0.5\textwidth}
    \begin{itemize}
      \item Análise de dados em arquivo CSV (esportes, clima, finanças)
      \item Calculadora científica / de unidades
      \item Quiz com banco de perguntas em arquivo
      \item Gerador de relatórios (texto + gráficos)
    \end{itemize}
  \end{columns}
\end{frame}

\begin{frame}{Sugestões por engenharia}
  \small
  \begin{tabular}{@{}p{3.6cm}p{8.6cm}@{}}
    \toprule
    \textbf{Civil / Produção} & Gestão de insumos e custos com dicionários e arquivos; controle de estoque; análise de produção e paradas (CSV + Pandas); orçamento de obra.\\
    \textbf{Elétrica / Mecânica} & Simulador de circuitos (leis de Kirchhoff com NumPy); tabelas termodinâmicas; análise de ensaios (ajuste e gráfico); dimensionamento de componentes.\\
    \textbf{Ambiental / Cartográfica} & Séries históricas de temperatura/chuva; conversor de coordenadas geográficas; cálculo de distâncias e áreas; mapas de calor em matriz.\\
    \textbf{Computação} & Jogos com matrizes (velha, campo minado); simulações; recursão e ordenação aplicadas; mini-sistema com módulos.\\
    \bottomrule
  \end{tabular}
\end{frame}

\begin{frame}{Dados para analisar}
  \begin{itemize}
    \item Arquivos de exemplo na pasta do curso: produção de peças (\texttt{dadosproducao.csv}) e produção/paradas mensais
    \item Dados abertos: IBGE, INMET (clima), Banco Mundial, Gapminder, Portal da Transparência
    \item \textbf{Seus próprios dados}: medições de laboratório, planilhas de estágio ou de iniciação científica (com autorização)
  \end{itemize}
  \begin{block}{Boa pergunta de projeto}
    ``Qual máquina/turno tem a maior taxa de defeitos?'' \(\cdot\) ``Como a temperatura média evoluiu em 30 anos?'' \(\cdot\) ``Qual a resistência de um componente a partir de ensaios?''\\
    \small\emph{Uma pergunta clara \(\Rightarrow\) um projeto focado.}
  \end{block}
\end{frame}

\begin{frame}{Hoje: formação dos grupos e esboço inicial}
  \begin{enumerate}
    \item \textbf{Forme o grupo} (2 ou 3 pessoas) e escolha um nome
    \item \textbf{Escolha o tema} (ou proponha o seu) e responda por escrito:
    \begin{itemize}\small
      \item Qual problema vamos resolver? Para quem?
      \item Quais as entradas e as saídas?
      \item Que dados usaremos? De onde vêm?
    \end{itemize}
    \item \textbf{Esboço inicial} (meia página): 3 a 5 funcionalidades principais
    \item Entregue o esboço ao professor ao final da aula
  \end{enumerate}
\end{frame}

\begin{frame}{Cronograma do projeto}
  \centering\small
  \begin{tabular}{cll}
    \toprule
    \textbf{Aula} & \textbf{Etapa} & \textbf{Entrega}\\
    \midrule
    26 & Lançamento & grupos e esboço\\
    27--29 & Matplotlib e Pandas (ferramentas) & exercícios de laboratório\\
    30 & Especificação & documento de especificação\\
    31 & Implementação 1 & esqueleto funcionando\\
    32 & Implementação 2 & arquivos e erros\\
    33 & Refatoração e testes & código limpo e testes\\
    34--35 & Apresentação & código + demonstração\\
    \bottomrule
  \end{tabular}
\end{frame}

\begin{frame}{Critérios de avaliação}
  \begin{columns}[T]
    \column{0.5\textwidth}
    \begin{block}{Funcionalidade}
      O programa atende aos requisitos especificados e funciona sem ``quebrar''?
    \end{block}
    \begin{block}{Organização do código}
      Bem estruturado, nomes claros, fácil de entender?
    \end{block}
    \column{0.5\textwidth}
    \begin{block}{Modularização}
      Código organizado em funções e módulos de forma lógica?
    \end{block}
    \begin{block}{Clareza na explicação}
      Os alunos explicam as decisões de projeto e a implementação?
    \end{block}
  \end{columns}
  \vspace{1ex}
  \small Todos os integrantes devem saber explicar \emph{qualquer} parte do código.
\end{frame}

\begin{frame}{Dicas para um bom começo}
  \begin{itemize}
    \item \textbf{Comece pequeno}: uma versão simples que funciona vale mais que uma ambiciosa que não roda
    \item Crie a pasta do projeto, um arquivo principal e \alert{use controle de versões} se conhecer Git
    \item Divida as tarefas, mas \textbf{integrem cedo e com frequência}
    \item Peça feedback ao professor em toda aula de implementação
    \item Código em equipe: combinem o estilo (nomes, indentação, docstrings)
  \end{itemize}
  \begin{exampleblock}{Próxima aula (30)}
    Especificação detalhada do problema, das funcionalidades e da estrutura do programa.
  \end{exampleblock}
\end{frame}

\end{document}
